
\section{Сравнение генерации задач по математике с помощью нейросети и с помощью
комбинированного алгоритма}

Появление генеративных нейросетей создало естественный интерес к автоматизации этого процесса: 
если ИИ способен генерировать тексты и код, то почему бы не поручить ему создание математических задач? 
Однако на практике применение нейросетей в образовательной среде сталкивается с рядом системных ограничений, 
касающихся качества, воспроизводимости и экономической целесообразности.

Для достижения поставленной цели было сформулированно четыре критерия оценки методов генерации:
\begin{itemize}
    \item уникальность вариантов,
    \item равномерность распределения параметров в допустимом диапазоне,
    \item вычислительная эффективность,
    \item возможность локального воспроизведения без зависимости от внешних API.
\end{itemize}
I.
В качестве тестовой задачи была выбрана задача типа №9299454 из открытого банка ЕГЭ по теме «Энергия конденсатора»,
 обладающая двумя степенями свободы — ёмкостью $C$ и напряжением $U$ в диапазоне 1–200. 
На всех изображениях далее по вертикали будет $C$ , а по горизонтали $U$ .

%Что нам дали нейросетевые подходы: эмпирические наблюдения и ограничения.
\subsection{Параллельная генерация}
Параллельная генерация на рис.~\ref{ris:pic-10} продемонстрировала наиболее низкие показатели: из 80 запрошенных вариантов было получено менее 15 уникальных значений 
параметра $C$. Наблюдается выраженная кластеризация: значения $C=25$ и
$C=50$ повторялись более $15$ и $20$ раз соответственно. Аналогичная картина характерна и для параметра 
$U$, где доминируют «круглые» числа ($100, 50, 20, 12$). Это свидетельствует о том, что стохастическая природа генерации и склонность модели к 
воспроизведению наиболее вероятных паттернов делают параллельный запрос непригодным для создания уникальных экзаменационных материалов.
\begin{figure}[h]
\vfill
\begin{minipage}[h]{0.8\linewidth}
\center{\includegraphics[width=1\linewidth]{pic-10.png}}
\end{minipage}
\caption{Параллельная генерация нейросетью ($100$ значений, диапазон $1-200$)}
\label{ris:pic-10}
\end{figure}


\subsection{Прямой запрос}
Прямой запрос на генерацию $40$ значений показал на рис.~\ref{ris:pic-11} систематические искажения распределения: нейросеть тяготеет к «удобным» числам 
($5, 10, 25, 50, 100, 200$), оставляя целые диапазоны (например, $26–39$, $51–79$) непокрытыми. Значения 
$U$ концентрируются в области $2–20$, хотя диапазон допускает значения до 200. Таким образом, даже при явном указании требований к равномерности,
модель не способна обеспечить статистически корректное покрытие параметрического пространства.
\begin{figure}[h]
\begin{minipage}[h]{0.8\linewidth}
\center{\includegraphics[width=1\linewidth]{pic-11.png}}
\end{minipage}
\caption{Прямой запрос нейросети($40$ значений, диапазон $1-200$)}
\label{ris:pic-11}
\end{figure}


\subsection{Последовательная генерация}
Последовательная генерация  на рис.~\ref{ris:pic-12} в рамках единого диалога показала наилучший результат среди чисто нейросетевых методов: значения параметров 
демонстрируют хороший разброс по всему диапазону, включая нетривиальные числа ($25, 128, 162, 192$ и др.). Однако и здесь сохраняются проблемы:
 отдельные повторения ($50, 100, 150$), неравномерная плотность распределения, а также линейный рост времени генерации с увеличением контекста.
Уже на $20–30$ итерации задержки становятся критическими, что ограничивает масштабируемость подхода.
Общий вывод по нейросетевым методам: генеративные модели, включая Qwen-3.6-Plus, ценны для прототипирования, адаптации формулировок и генерации 
кода, но их применение в качестве основного инструмента для пакетной генерации экзаменационных вариантов нецелесообразно из-за стохастичности выдачи, 
кластеризации на «круглых» числах и высоких вычислительных издержек.
\begin{figure}[h]
\begin{minipage}[h]{0.8\linewidth}
\center{\includegraphics[width=1\linewidth]{pic-12.png}}
\end{minipage}
\caption{Последовательная генерация нейросетью ($40$ значений, диапазон $1-200$)}
\label{ris:pic-12}
\end{figure}


\subsection{Гибридный подход}
Гибридный подход на рис.~\ref{ris:pic-13}, при котором нейросеть используется для генерации детерминированного JavaScript-кода, 
а сам код — для производства задач, продемонстрировал хорошее покрытие параметрического пространства
на диапазоне $1-200$.Заметна высокая уникальность значений, распределение
визуально равномерное без явных кластеров. Отсутствует тяготение к круглым
числам на рис.~\ref{ris:pic-13}.  Визуальный анализ тепловой карты распределения
$200$ сгенерированных пар ($C,U$) показал практически равномерное покрытие параметрического пространства $1–200$ без выраженных кластеров и 
«пустых зон».Для формальной оценки равномерности мы разбили пространство на $36$ ячеек и вычислили отклонения наблюдаемых частот от ожидаемых 
при гипотетически равномерном распределении. Диапазон отклонений составил от $−3$ до $+5$ точек на ячейку, что менее $3$ процента от объёма выборки.
Красные и синие области на карте отклонений распределены хаотично, без образования регулярных паттернов, что подтверждает отсутствие систематических смещений.
Статистическая проверка с помощью теста дала p-value = 0.89 > 0.05, что означает: наблюдаемое отклонение от равномерного распределения статистически незначимо. В отличие от нейросетевых методов, 
где кластеризация на круглых числах проявилась бы в виде выраженных «горячих пятен», детерминированная генерация обеспечивает равномерное заполнение всего пространства.
\begin{figure}[h]
    \centering
    \begin{minipage}[h]{0.8\linewidth}
        \centering
        \includegraphics[width=1\linewidth]{pic-13.png}
        \\\textit{''Час ЕГЭ'' ($200$ значений, диапазон $1-200$)}
    \end{minipage}
    
    \vspace{0.5cm}
    
    \begin{minipage}[h]{0.8\linewidth}
        \centering
        \includegraphics[width=1\linewidth]{pic-14.png}
        \\\textit{Описание графиков распределения параметров ``Час ЕГЭ'', диапазон $1-200$}
    \end{minipage}
    
    \caption{Распределение параметров при генерации задач: тепловая карта и отклонения от равномерности}
    \label{ris:pic-13}
\end{figure}

Таким образом, генеративные нейросети являются ценным инструментом для прототипирования и адаптации формулировок,
однако для задач пакетной генерации экзаменационных вариантов с требованиями к уникальности, равномерности и масштабируемости предпочтителен гибридный подход,
сочетающий гибкость ИИ и предсказуемость детерминированных алгоритмов.


%\begin{figure}[h]
%\begin{minipage}[h]{1\linewidth}
%\center{
\includegraphics[width=1\linewidth]{506300-sdam-2.png}} \\а)
%\end{minipage}
%\vfill
%\begin{minipage}[h]{0.95\linewidth}
%\center{%
%    \includegraphics[width=1\linewidth]{506300-4-1.png}
%    \\
%    \includegraphics[width=1\linewidth]{506300-4-2.png}
%} б) \\
%\end{minipage}
%\caption{Задачи, соответствующие шаблону 506300:
%\\
%а) исходная; б) сгенерированные}
%\label{ris:506300}
%\end{figure}
